\documentclass[12pt]{article}

% support for images
\usepackage{graphicx}

% font
\usepackage{tgheros}
\renewcommand*\familydefault{\sfdefault}
\usepackage[T1]{fontenc}


% justify
\usepackage{ragged2e}

% indents
\usepackage{indentfirst}
\setlength{\parindent}{0pt}
\setlength{\parskip}{0pt}
\setlength{\JustifyingParindent}{1.25cm}
% source: https://tex.stackexchange.com/a/121230
\renewenvironment{justify}{%
	\trivlist
	\justifying
	\itemindent\JustifyingParindent
	\item\relax
	}{%
	\endtrivlist
}
\justifying

% fonts, margins, etc...
\usepackage{setspace}
\usepackage{authoraftertitle}
\usepackage{tikz}
\usepackage{anyfontsize}
\usetikzlibrary{positioning,calc}

% language
\usepackage{polski}
\usepackage[utf8]{inputenc}
\usepackage[english,polish]{babel}

% bibliography
\usepackage{biblatex}
\addbibresource{../bibliography.bib}

% configure code fields
\usepackage{listings}
\lstdefinestyle{customcpp}{
  belowcaptionskip=1\baselineskip,
  breaklines=true,
  frame=L,
  xleftmargin=\parindent,
  language=C++,
  showstringspaces=false,
  basicstyle=\fontsize{9pt}{9}\ttfamily,
  keywordstyle=\bfseries\color{green!40!black},
  commentstyle=\itshape\color{purple!40!black},
  stringstyle=\color{orange},
}
\lstset{
	backgroundcolor=\color{white},
	language=C++,
	style=customcpp,
	tabsize=2
}

% set margins
\usepackage[
	a4paper,
	left=2.5cm,
	right=2.5cm,
	top=2.5cm,
	bottom=2.5cm,
	bindingoffset=1cm,
	headheight=1.25cm,
	footheight=1.35cm
]{geometry}

% configure header and footer (with page numbering)
\usepackage[lastpage,user]{zref}
\usepackage{fancyhdr}
\renewcommand{\headrulewidth}{0pt}
\fancyhf{}
\pagestyle{fancy}
\cfoot{\thepage\ z \zpageref{LastPage}}

% document info
\title{Implementacja i badanie wydajności serializacji wykonanej w paradygmacie meta-programowania względem standardowych rozwiązań w języku C++}
\author{inż. Krzysztof MOCHOCKI}
\date{}

% configure enumerator
\renewcommand{\labelenumi}{\arabic{enumi}.}
\renewcommand{\labelenumii}{\arabic{enumi}.\arabic{enumii}.}
\renewcommand{\labelenumiii}{\arabic{enumi}.\arabic{enumii}.\arabic{enumiii}.}
\renewcommand{\labelenumiv}{\arabic{enumi}.\arabic{enumii}.\arabic{enumiii}.\arabic{enumiv}.}

% configure sections
\usepackage{titlesec}
\titleformat{\section}{\normalfont\Large\bfseries}{\thesection.}{1em}{}
\titleformat{\subsection}{\normalfont\large\bfseries}{\thesubsection.}{1em}{}
\titleformat{\subsubsection}{\normalfont\normalsize\bfseries}{\thesubsubsection.}{1em}{}

% function for disabling hyphenation
\newcommand{\disablewordwrap}{%
	\setlength{\hsize}{0.9\hsize}
	% \tolerance=1
	\hyphenpenalty=10000
	\hbadness=10000}

% configuration of table of contents
\usepackage{tocloft}
\setcounter{tocdepth}{2}
\renewcommand{\cftsecleader}{\cftdotfill{\cftdotsep}}
\usepackage[hidelinks]{hyperref}

% add page break with new section
\newcommand{\sectionbreak}{\pagebreak}

% configure quotation
\usepackage{etoolbox}
\AtBeginEnvironment{quotation}{\itshape}

% configuration of figure captions
% source: https://tex.stackexchange.com/a/276796
\usepackage{caption}
\captionsetup[figure]{labelfont={bf}, labelformat={default}, labelsep=period, name={Wykres}}

% configure placing of figures in-text
% source: https://tex.stackexchange.com/a/8633
\usepackage{float}

% line spread
\linespread{1.5}

% aliases
\newcommand{\n}{\newline}

\begin{document}

	% make default captions disappear of list of tables and list of figures
	% to set custom, numbered mentioned in ToC
	\renewcommand*\listfigurename{}
	\renewcommand*\listoftables{}

	% tittle page
	\begin{titlepage}
		\clearpage
		\topskip0pt
		\centering

		\includegraphics[width=5cm, keepaspectratio=true]{./img/black_and_white_polsl_logo.png}

		{\LARGE\bfseries\ PRACA MAGISTERSKA}

		\vspace*{1cm}

		{\LARGE \MyTitle}

		\Large\bfseries\

		\begin{spacing}{1.125}
			\MyAuthor\

			000000000
			\vspace*{1cm}

			Kierunek: Informatyka Przemysłowa

			Specjalność: Inteligentne Systemy Przemysłowe

			\vspace*{1cm}

			PROMOTOR

			dr. inż. Adrian SMAGÓR

			\vspace*{0.5cm}

			KATEDRA Informatyki Przemysłowej

			Wydział Inżynierii Materiałowej

			\vspace*{\vfill}

			KATOWICE 2023
		\end{spacing}

		\thispagestyle{empty}
	\end{titlepage}

	% detail page
	{
		\disablewordwrap
		\large

		\vspace*{\vfill}

		{\bfseries\ Tytuł pracy:} \MyTitle

		\vspace*{\vfill}

		{\bfseries\ Streszczenie:} <tu będzie streszczenie>

		\vspace*{\vfill}

		{\bfseries\ Słowa kluczowe:} serializacja, deserializacja, C++, szablony, wydajność, meta-programowanie

		\vspace*{\vfill}

		{\bfseries\ Thesis title:} Implementation and performance testing of serialization performed in the meta-programming paradigm against standard solutions in the C++ language

		\vspace*{\vfill}

		{\bfseries\ Abstract:} <tu będzie streszczenie po angielskiemu>

		\vspace*{\vfill}

		{\bfseries\ Keywords:} serialization, deserialization, C++, templates, performance, meta-programming

		\vspace*{\vfill}

		\thispagestyle{empty}
		\newpage
	}

	% blank page
	{
		 \vspace*{\vfill}
		\thispagestyle{empty}
		\newpage
	}

	% table of contents
	{
		\tableofcontents
		\thispagestyle{fancy}
		\newpage
	}

	% introduction
	{
		\section{Wstęp}

		Serializacja w językach wysokiego poziomu, takich jak JavaScript czy Python jest, w przypadku standardowych formatów
		reprezentacji danych, echem przeszłości, ponieważ języki te posiadają natywne wsparcie do reprezentacji struktur danych
		zapewnione przez ichniejsze biblioteki standardowe.\n

		Niestety w momencie powstawania bieżącej pracy programiści języka C++ wciąż muszą uwzględniać problematykę
		serializacji podczas projektowania aplikacji oraz doboru bibliotek. Standard języka C++ w roku 2023 posiada szeroką
		paletę narzędzi do pisania szablonów klas i struktur, niestety nie posiada żadnych narzędzi do łatwego przeglądania
		struktur danych. Jest to oczywiście podyktowane względami wydajnościowymi oraz chęcią zachowania wstecznej
		kompatybilnością z językiem C. Istnieją rozszerzenia (na przykład reflexpr \cite*{cpp_extension_reflexpr_iso}), które
		umożliwiają iterowanie po elementach struktur danych, czy ich serializację, jednakże nie są one wciąż wdrożone do
		standardu \cite*{cpp_extension_reflexpr_short}.\n

		\begin{figure}[H]
			\centering
			\includegraphics{./charts/output_with_charts_as_images/json_vs_other_formats.png}
			\caption{Popularność haseł według strony Google \cite*{google_trends_json_api_vs_rest_of_the_world}}
			\medskip
		\end{figure}

		Niniejsza praca skupia się wyłącznie na formacie JSON \cite*{json_iso}, ponieważ jak można zauważyć na powyższym wykresie, jest on
		najbardziej powszechnym standardem komunikacji na świecie.\n

		Aktualnie możemy rozróżnić trzy podstawowe podejścia do tego problemu. Pierwsze, gdzie programista tylko wspomaga się
		zewnętrzną biblioteką, której jedynym zadaniem jest zapewnienie zgodności ze standardem formatu, zarówno podczas serializacji
		oraz deserializacji, natomiast nie zapewnia refleksji typów złożonych, których serializację bądź deserializację musi
		dokonać samodzielnie.\n

		Kolejnym podejściem do problemu jest rozwiązanie oparte na bibliotece, która również jak powyżej zawiera funkcję
		serializacji i deserializacji typów prostych, jednakże dodatkowo zawiera funkcję zewnętrznej refleksji typów,
		która najczęściej już na podstawie statycznych referencji do pola klasy, jest wstanie wydedukować typy i automatycznie
		poprowadzić proces serializacji lub deserializacji.\n

		Trzecim podejściem jest zastosowanie wewnętrznej refleksji. W innych językach najczęściej jest dostarczona przez język.
		Przykładowo w języku Python jest ona wyrażona w postaci pola {\it \_\_dict\_\_}, które dostarcza programiście całą klasę
		w postaci dynamicznej struktury danych typu klucz-wartość: słownika, gdzie kluczem jest nazwa pola klasy o typie
		łańcucha znakowego, natomiast wartością jest konkretna wartość w obecnej instancji klasy. Z punktu widzenia tematu
		niniejszej pracy, daje to pełny i prosty wgląd w strukturę zarówno klasy jak i hierarchii klas.
	}

	% theory
	{
		\section{Teoria}

		{
			\subsection{Definicja serializacji oraz deserializacji}

			Jak podaje portal {\it hazelcast.com}\cite{definitions_serial_and_deserial} definicja serializacji oraz deserializacji brzmi:

			\begin{quotation}
				{\bf Serialization} is the process of converting a data object - a combination of code
				and data represented within a region of data storage - into a series of bytes that
				saves the state of the object in an easily transmittable form. [\ldots]\n

				The reverse process - constructing a data structure or object from a series of
				bytes - is {\bf deserialization}. The deserialization process recreates the object,
				thus making the data easier to read and modify as a native structure in a programming language.
			\end{quotation}

			{\noindent Tłumacząc na język polski (tłumaczenie własne):}

			\begin{quotation}
				{\bf Serializacja} to proces konwersji obiektu danych - kombinacji kodu i danych
				reprezentowanych w regionie przechowywania danych - w serię bajtów zapisujących
				stan obiektu w łatwej do przesłania formie. [\ldots] \n

				Proces odwrotny - konstruowanie struktury danych lub obiektu z~szeregu
				bajtów - to {\bf deserializacja}. Proces deserializacji odtwarza obiekt,
				dzięki temu dane są łatwiejsze do odczytania i modyfikacji jako
				natywna struktura w języku programowania.
			\end{quotation}
		}

		{
			\newpage
			\subsection{Format JSON}

			Jest lekkim formatem danych wywodzącym się z języka JavaScript. Jak pokazano na wykresie numer jeden,
			JSON jest formatem o największej rozpoznawalności wśród tekstowych reprezentacji danych. Swój sukces
			zawdzięcza swojej formie: bardzo czytelnej dla człowieka, jednocześnie wydajnej dla komputerów zarówno
			podczas przetwarzania jak i generowania. Cecha ta sprawia, że praktycznie każdy współcześnie używany język
			programowania ma jego obsługę\cite{json_short} w~bibliotece standardowej lub posiada zastęp zewnętrznych
			bibliotek.\n

			Poddając JSON usystematyzowaniu, jest to format danych o strukturze drzewiastej, który technicznie
			może być zagnieżdżony w stopniu nieskończonym, w przeciwieństwie do formatów płaskich, takich jak CSV.\n
		}

		{
			\subsection{Definicja metaprogramowania}

			Opis, a zarazem definicja, czym metaprogramowanie jest zostało opisane niniejszym cytatem \cite{metaprogramming_definition}:

			\begin{quotation}
					In linguistics, a metalanguage is defined as follows [Dict]:\n

				“any language or symbolic system used to discuss, describe, or analyze another
					language or symbolic system”\n

					This definition characterizes the main idea of metaprogramming, which involves writing
				programs related by the meta-relationship, i.e. the relationship of “being about”. A program that
				manipulates another program is clearly an instance of metaprogramming. [\ldots]
			\end{quotation}

			\newpage
			{\noindent Tłumacząc na język polski (tłumaczenie własne):}

			\begin{quotation}
					W językoznawstwie metajęzyk definiuje się w następujący sposób [Słownik]:\n

				„dowolny język lub system symboliczny używany do omawiania, opisywania lub analizowania
					innego języka lub systemu symbolicznego”\n

					Ta definicja charakteryzuje główną ideę metaprogramowania, która obejmuje pisanie
				programów powiązanych meta-relacją, czyli relacją „bycia”. Program, który
				manipuluje innym programem jest ewidentnie przykładem metaprogramowania. [\ldots]
			\end{quotation}

			Paradygmat meta programowania w przypadku niniejszej pracy będzie można zaobserwować przy
			rozwiązaniach z zastosowaniem szablonów, a także nowinki językowej w języku C++: konceptów.
			Meta-relacja będzie się ujawniać, gdy kompilator (który w tym wypadku będzie wykonawcą meta języka),
			będzie wykonywał i generował automatycznie procedury refleksji typów, a także procedury ich serializacji
			oraz deserializacji.\n
		}

		{
			\newpage
			\subsection{Definicja programowania refleksyjnego}

			Definicję programowania zorientowanego na refleksję podaje Jonathan M. Sobel oraz Daniel P. Friedman
			w swoim materiale {\it An Introduction to Refection-Oriented Programming}\cite{reflection_definition},
			gdzie możemy przeczytać:

			\begin{quotation}
				[\ldots] reflective computational systems allow computations to observe and modify properties
				of their own behavior, especially properties that are typically observed only from some
				external, meta-level view point. [\ldots]
			\end{quotation}

			{\noindent Tłumacząc na język polski (tłumaczenie własne):}

			\begin{quotation}
				[\ldots] refleksyjne systemy obliczeniowe umożliwiają przetwarzanie poprzez obserwację i~modyfikację
				właściwości własnego zachowania, zwłaszcza właściwości, które zazwyczaj obserwuje się tylko z~zewnętrznego,
				meta-poziomowego punktu widzenia. [\ldots]
			\end{quotation}

			Dobrym przykładem programowania refleksyjnego, jest kompletna kontrola nad typami danych w języku Python,
			gdzie istnieje pełna możliwość dynamicznego przeglądania, dodawania, czy usuwania pól w klasie (lub wyłącznie w jednym obiekcie)
			podczas działania programu. Jest to na tyle rozbudowany obszar w~tym języku, że istnieją biblioteki takie jak
			{\it dataclasses}\cite{python_dataclass}, które podmieniają prostą definicję klasy z wyłącznie statycznymi polami,
			na definicję klasy z stworzonymi konstruktorami, reprezentacjami i zwykłymi polami, należącymi do obiektu.

			Warto nadmienić, że język C++ domyślnie nie wspiera programowania refleksyjnego, co jest przyczółkiem problematyki
			całej niniejszej pracy.
		}
	}

	% \subsection{Pod Wstęp}
	% \subsubsection{Pod Pod Wstęp}

	% \begin{enumerate}
	% 	\item One
	% 	\item Two
	% 	\item Three
	% 	\begin{enumerate}
	% 		\item Three point one
	% 		\begin{enumerate}
	% 		\item Three point one, point one
	% 			\begin{enumerate}
	% 			\item Three point one, point one, point one 1234567890
	% 			\item Three point one, point one, point two
	% 			\end{enumerate}
	% 		\end{enumerate}
	% 	\end{enumerate}
	% 	\item Four
	% 	\item Five
	% 	\end{enumerate}

	% bibliography
	{
		\newpage
		\section{Bibliografia}
		\printbibliography[heading=none]
	}

	% list of tables
	{
		\newpage
		\section{Spis Tabel}
		\listoftables
		\thispagestyle{fancy}
		\newpage
	}

	% list of figures
	{
		\newpage
		\section{Spis Wykresów}
		\listoffigures
		\thispagestyle{fancy}
		\newpage
	}

\end{document}
